\section{Спецификация VAR/SVAR}

Модель Sims и Zha строится на следующих стилизованных фактах:

\begin{enumerate}
\item При использовании инноваций денежной массы для идентификации монетарных шоков возникает «загадка ликвидности» (liquidity puzzle): денежное сжатие не приводит к росту процентной ставки.

\item При использовании инноваций краткосрочной процентной ставки возникает «загадка цен» (price puzzle): денежное сжатие не приводит к снижению цен.

\item Наблюдаемая связь между высокой процентной ставкой и последующим снижением выпуска не обязательно означает, что жёсткая денежно-кредитная политика вызывает спад. Высокая ставка может отражать реакцию центрального банка на уже существующее инфляционное давление.

\item Большая часть изменений денежной массы и процентной ставки представляет собой систематическую реакцию денежно-кредитной политики на состояние экономики, а не экзогенные монетарные шоки.

\item Цены сырья содержат информацию о будущей инфляции и могут использоваться для отделения собственно монетарных шоков от реакции политики на изменение экономической конъюнктуры.

\end{enumerate}

Модель также опирается на следующие предпосылки:

\begin{enumerate}
\item В экономике существует инерционный сектор, переменные которого реагируют на денежно-кредитную политику с лагом.

\item Не все переменные частного сектора являются инерционными. Аукционные рыночные цены, прежде всего цены сырья, способны реагировать на монетарные шоки внутри периода.

\item Центральный банк наблюдает цены сырья и может использовать их при принятии решений по денежно-кредитной политике.

\item Модель должна отделять систематическую реакцию денежно-кредитной политики от собственно структурного монетарного шока.

\end{enumerate}

Теоретическая мотивация этих предпосылок даётся через DSGE-модель с инерционной подстройкой. В ней есть переменные, которые меняются медленно из-за издержек корректировки, но есть и более гибкие цены. Денежная политика задаётся правилом для краткосрочной ставки, которое реагирует на уровень цен и цены сырья, а также содержит отдельный монетарный шок. Экзогенные шоки моделируются как авторегрессионные процессы с белыми шумами.

В исходной статье модель оценивается на квартальных данных по экономике США за 1964--1994 годы. Все ряды, кроме показателя банкротств, взяты из базы Citibase; данные по банкротствам были отдельно предоставлены авторам и включают личные и корпоративные банкротства.

Эмпирически авторы используют идентифицированную структурную VAR-модель. Вектор переменных включает цены сырья, денежный агрегат, процентную ставку, цены промежуточных материалов, дефлятор выпуска, заработную плату, реальный выпуск и банкротства:


$$
x_t =
(P_{cm,t}, M_t, R_t, P_{im,t}, P_{y,t}, W_t, y_t, T_{bk,t})'.
$$

Структурная форма модели записывается как

$$
\Gamma(L)x_t=\varepsilon_t,
$$

Приведённая форма имеет вид

$$
x_t =
B_1 x_{t-1}
+\dots+
B_p x_{t-p}
+
u_t.
$$

Связь между остатками приведённой VAR и структурными шоками задаётся как

$$
\varepsilon_t=\Gamma_0u_t,
\qquad
\Gamma_0\Sigma_u\Gamma_0'=\Sigma_\varepsilon.
$$

что можно записать более развернуто:
$$
\Gamma_0 x_t
=
-\Gamma_1 x_{t-1}
-\dots
-\Gamma_p x_{t-p}
+
\varepsilon_t.
$$

Авторы нормируют структурные шоки так, что

$$
\Sigma_\varepsilon=I.
$$

Приведённая VAR оценивается методом наименьших квадратов. После этого на матрицу текущих связей $\Gamma_0$ накладываются идентифицирующие ограничения. Они позволяют выделить шоки денежного спроса, денежной политики и частного сектора. Сама структурная матрица $\Gamma_0$ оценивается методом максимального правдоподобия при заданных ограничениях.

Эти ограничения задают структуру матрицы $\Gamma_0$ и определяют систему одновременных уравнений для текущих значений переменных. Через неё формализуются предположения о том, какие переменные могут реагировать друг на друга внутри периода, а какие реагируют только с лагом.

Этой матрице соответствует система текущих уравнений:

$$
\left\{
\begin{array}{ll}
P_{cm}:&
\gamma_{11}P_{cm,t}+\gamma_{12}M_t+\gamma_{13}R_t+\gamma_{14}P_{im,t}
+\gamma_{15}P_{y,t}+\gamma_{16}W_t+\gamma_{17}y_t+\gamma_{18}T_{bk,t}
=
\text{лаги}+\varepsilon^{Pcm}_t,
\\[0.5em]

MD:&
\gamma_{22}M_t+\gamma_{23}R_t-\gamma_{22}P_{y,t}-\gamma_{22}y_t
=
\text{лаги}+\varepsilon^{MD}_t,
\\[0.5em]

MS:&
\gamma_{31}P_{cm,t}+\gamma_{32}M_t+\gamma_{33}R_t+\gamma_{38}T_{bk,t}
=
\text{лаги}+\varepsilon^{MS}_t,
\\[0.5em]

P_{im}:&
\gamma_{41}P_{cm,t}+\gamma_{44}P_{im,t}+\gamma_{45}P_{y,t}
+\gamma_{46}W_t+\gamma_{47}y_t+\gamma_{48}T_{bk,t}
=
\text{лаги}+\varepsilon^{Pim}_t,
\\[0.5em]

P_y:&
\gamma_{51}P_{cm,t}+\gamma_{55}P_{y,t}+\gamma_{56}W_t
+\gamma_{57}y_t+\gamma_{58}T_{bk,t}
=
\text{лаги}+\varepsilon^{Py}_t,
\\[0.5em]

W:&
\gamma_{61}P_{cm,t}+\gamma_{66}W_t+\gamma_{67}y_t+\gamma_{68}T_{bk,t}
=
\text{лаги}+\varepsilon^{W}_t,
\\[0.5em]

y:&
\gamma_{71}P_{cm,t}+\gamma_{77}y_t+\gamma_{78}T_{bk,t}
=
\text{лаги}+\varepsilon^{y}_t,
\\[0.5em]

T_{bk}:&
\gamma_{81}P_{cm,t}+\gamma_{88}T_{bk,t}
=
\text{лаги}+\varepsilon^{Tbk}_t.
\end{array}
\right.
$$

В TR-модели структура аналогична, но из уравнения денежной политики исключается $T_{bk,t}$:

$$
MS:
\quad
\gamma_{31}P_{cm,t}
+\gamma_{32}M_t
+\gamma_{33}R_t
=
\text{лаги}
+
\varepsilon^{MS}_t.
$$

Экономическая логика ограничений следующая. Уравнение $MD$ задаёт денежный спрос: денежный агрегат связан со ставкой, ценами и выпуском. Уравнение $MS$ задаёт денежную политику: политика может внутри квартала реагировать на деньги и быстро наблюдаемые сырьевые цены, а в M2-модели также на банкротства. Остальные уравнения образуют медленный блок частного сектора: $P_{im}$, $P_y$, $W$, $y$ и $T_{bk}$ не реагируют внутри квартала на $M$ и $R$, но могут реагировать на $P_{cm}$.

В статье оцениваются две основные версии модели. В первой денежный блок представлен агрегатом M2 и 90-дневной ставкой по казначейским векселям. Во второй используются совокупные резервы банков и ставка федеральных фондов. Это позволяет проверить, насколько выводы о роли денежно-кредитных шоков зависят от выбора денежного индикатора и процентной ставки.
